\section*{Chapter \ref{ch:s1:volatility-targeting-and-risk-managed-portfolios} --- Volatility Targeting and Risk-Managed Portfolios}

\begin{solution}{exo:s1:volatility-targeting-and-risk-managed-portfolios:1}
$0.10/0.12 = 0.83$ and $0.10/0.40 = 0.25$.
\end{solution}

\begin{solution}{exo:s1:volatility-targeting-and-risk-managed-portfolios:2}
$0.002 \times (0.73 - 0.41)/0.004 = 0.16$: about 16\% of the day's volume (the simulation's 14.5\% nets the day's price drift in the funds' exposure).
\end{solution}

\begin{solution}{exo:s1:volatility-targeting-and-risk-managed-portfolios:3}
Exposure moves with the square of the forecast's ratio: a forecast 10\% higher cuts exposure by about 17\% instead of 9\%, so every forecast change is a
larger trade.
\end{solution}

\begin{solution}{exo:s1:volatility-targeting-and-risk-managed-portfolios:4}
$0.06/0.16 = 0.375$. A ten-year Sharpe ratio has a standard error of about $1/\sqrt{10} = 0.32$ (the samples range from $-0.44$ to 1.18). The mean gain over
twenty samples has a standard error of $0.22/\sqrt{20} = 0.049$, so 0.028 is not significant even with twenty decades.
\end{solution}

\begin{solution}{exo:s1:volatility-targeting-and-risk-managed-portfolios:5}
In equities volatility rises when prices fall (the leverage effect), so volatility is high after losses and cutting exposure then avoids the worst
continuation; bonds lack that link, so their volatility says little about their next returns and scaling adds only noise.
\end{solution}

\begin{solution}{exo:s1:volatility-targeting-and-risk-managed-portfolios:6}
A volatility spike makes funds sell; their selling moves prices further; the larger moves raise the funds' volatility forecasts; the higher forecasts
make them sell more. Impact and the forecast's memory decide how far the loop runs.
\end{solution}

\begin{solution}{exo:s1:volatility-targeting-and-risk-managed-portfolios:7}
Before the spike the market's realised volatility rises from 13.0\% (no impact) to 18.7\%, and the funds' exposure falls from 0.73 to 0.46: their own
daily rebalancing, about 1.9\% of volume on an average day, moves prices, which raises the volatility they forecast, which makes them trade more. Large
enough, targeting funds generate part of the volatility they target.
\end{solution}

\begin{solution}{exo:s1:volatility-targeting-and-risk-managed-portfolios:8}
The spanning regression's implied strategy combines the managed and unmanaged portfolios with weights estimated on the whole sample; it is not
investable in real time. Compare the managed portfolio directly with the unmanaged one, out of sample, with costs, as Cederburg and co-authors did.
\end{solution}

\begin{solution}{pb:s1:volatility-targeting-and-risk-managed-portfolios:1}
\begin{enumerate}
\item A portfolio scaled by the inverse of forecast volatility; the net trading of funds following mechanical exposure rules.
\item A variance forecast, a scaling rule, a leverage cap and a rebalancing band.
\item Changes in volatility are not offset by proportional changes in expected returns, so taking less risk when volatility is high raises Sharpe ratios.
\item Across 103 strategies, managed portfolios do not systematically beat unmanaged ones out of sample.
\item Twenty ten-year samples of a GJR-GARCH market with a constant premium, constant or managed by last month's realised variance.
\item 0.49, 0.52 (inverse volatility) and 0.48 (inverse variance); 10 and 8 of 20 samples improved.
\item The gain (0.03) is a tenth of a decade's standard error (0.32).
\item The gain holds only for risk assets with a leverage effect; the synthetic classes, without one, change in both directions.
\item Funds with 0.2\% of market value targeting 10\%, a market trading 0.4\% a day, a 6\% fall and twenty days at 40\% volatility; all assumed.
\item 14.5\% of the next day's volume; 20.5\% of a day's volume over twenty days.
\item It deepens the twenty days' fall from 16.6\% to 18.7\% (impact 0.1) or 22.0\% (0.3).
\item It follows from the market's own recent volatility and the funds' public rules.
\item \textbf{Named result.} Managing the synthetic market by its volatility raised the Sharpe ratio by 0.028 on average over twenty samples (better in
10); after a 6\% fall, targeting funds holding 0.2\% of the market sold 14.5\% of the next day's volume.
\item It shrinks the tails and steadies the risk.
\item Real-time forecasts and constants, direct comparison with the unmanaged portfolio, costs, many periods.
\item Trading against deleveraging flow.
\item Chapter~19 targeted a trend book's volatility and found it levered up when the edge was small; here the question is the same for a market.
\item The funds' assets, rules and forecast windows, and the market's volume and depth at the close.
\item Their selling feeds their own forecasts.
\item A steadier risk, bought with selling into falls, which the market pays for in deeper falls.
\end{enumerate}
\end{solution}

\begin{solution}{iq:s1:volatility-targeting-and-risk-managed-portfolios:1}
If expected returns do not rise with volatility, a volatile period carries more risk per unit of expected return; cutting exposure then removes risk
faster than return, and the ratio over time improves.
\end{solution}

\begin{solution}{iq:s1:volatility-targeting-and-risk-managed-portfolios:2}
The alpha is earned by a combination of the managed and unmanaged portfolios with weights fitted on the whole sample, not available in real time;
out of sample the combination's weights are unstable and the advantage usually disappears.
\end{solution}

\begin{solution}{iq:s1:volatility-targeting-and-risk-managed-portfolios:3}
Volatility-targeting funds, risk parity, leveraged ETFs resetting at the close, and option dealers short gamma; estimate each group's assets, its
rule and its forecast window, and apply today's move to get tomorrow's trade.
\end{solution}

\begin{solution}{iq:s1:volatility-targeting-and-risk-managed-portfolios:4}
It reduces exposure to high-volatility periods and the severity of tail losses; it adds turnover, sells after falls, levers up in calm periods before
shocks, and ties the portfolio to crowded flows.
\end{solution}

\begin{solution}{iq:s1:volatility-targeting-and-risk-managed-portfolios:5}
Collect closing prices, update the forecast, compute the target exposure with caps and bands, check limits and liquidity, send orders near the close
with a fallback if markets are disrupted, and log every step for reconciliation.
\end{solution}

\begin{solution}{iq:s1:volatility-targeting-and-risk-managed-portfolios:6}
With $w_t = c/\sigma_t$ the return is $c\mu/\sigma_t + c\varepsilon_t$, with mean $c\mu E[1/\sigma]$ and variance $c^2(1 + \mu^2\operatorname{Var}(1/\sigma))$.
The constant position has Sharpe ratio $\mu/\sqrt{E[\sigma^2]}$. For small $\mu$ the ratio of the two is $E[1/\sigma]\sqrt{E[\sigma^2]} \ge E[1/\sigma]E[\sigma]
\ge 1$ by Jensen's inequality, with equality only when $\sigma_t$ is constant.
\end{solution}
